\section{局限}
\label{sec:limitations}

我们明确说明 \wcabench{} 无法支撑什么，因为基准的可信度取决于其自我设限的诚实程度。

\paragraph{外推风险（\task{4}）。}人类极限估计没有可观测真值，且天生易受结构性变化影响。历史数据无法预见未来由新硬件（如磁力魔方）、新训练体系或规则变化引发的跃变。变点模型（$3.02$~秒）、高斯过程叠加极值理论（$3.29$~秒）与指数衰减（$3.74$~秒）给出的三阶估计，与社区锚点（$\approx 2.37$~秒）的差异应被读作警示而非结论。本文报告的极限是模型条件性陈述，而非关于人类能力的事实。

\paragraph{地区不均衡。}各国家与大洲获得认证比赛的机会不均。由于参赛次数、对手强度等特征部分是机会的代理，跨地区泛化结论天然较弱。我们为透明起见报告大洲分层，但小分层（如大洋洲、南美）测试实例过少，无法得出可靠结论，且若干分层格子无定义。

\paragraph{自选择与混淆（\task{5}）。}参加某项目是与能力相关的选择。因此基于相关的迁移估计会高估效应，即便 DID 代理也依赖自选择可能违反的平行趋势假设。“可识别项目对数量”衡量的是数据对估计的支撑程度，而非因果效应的保证。在没有显式识别论证时，\wcabench{} 不应被用于对训练做出因果断言。

\paragraph{时间非平稳性。}训练窗口（2003--2022）与测试窗口（2025--2026）在硬件、规则与参与者人口构成上都不同。模型的退化可能反映真实分布漂移而非建模错误；\task{2} 的 Kendall $\kendall$ 随测试切片下降（$0.90 \to 0.64$）正与此一致。

\paragraph{小样本模式结果。}第~\ref{sec:results} 节的头条来自在单机上、单块 GPU 上运行的抽样测试窗口（第~\ref{sec:compute} 节）。当完整滚动窗口协议在全部 $1{,}719{,}290$ 条测试记录上执行时，绝对值会变化，排名并列（如 \task{2} 的 $0.7673$ 并列）也可能改变；三种子表（第~\ref{sec:results-multiseed} 节）正是为揭示这种抽样方差而纳入，但不能替代全量运行。我们报告抽样数字，是因为它们真实可核验，而非因为它们终局；已发布管线支持全量运行。

\paragraph{不平衡与稀疏类别。}\DNF{} 预测与技能迁移估计都受严重不平衡与稀疏影响。即便以进入决策点的轮次为条件，正类率也从部分项目的个位数到三阶盲拧的约 $95\%$ 不等，且许多项目对共同参赛选手过少，难以稳定估计。

\paragraph{原生后端与算力可复现性。}本次发布中，树模型基线运行在原生 XGBoost 上，GNN 运行在其原生 CUDA 路径上（第~\ref{sec:compute} 节）；每个报告都记录实际使用的设备、墙钟时间与 CPU/GPU 小时数。复现基线的读者应固定库版本与设备：不同的后端（例如树模型回退到 scikit-learn，或纯 CPU 图内核）无法精确复现这些数字，此前的回退后端结果与原生结果差异明显。

\paragraph{单一快照数据。}即便固定导出（v2.0.2），\wca{} 仍偶尔修正历史成绩，且数据库持续增长：每次每周导出都带新的顺序版本号并包含更多比赛。我们固定 2026-09-21 快照并记录之，因此确切的规模数字（如 $6{,}909{,}454$ 条成绩）是快照特定的，较新的导出会略多。将我们的计数与当前官方下载对比的读者，应把它们视为冻结基线而非实时值。

\section{伦理与更广泛影响}
\label{sec:ethics}

\paragraph{允许用途。}\wcabench{} 面向科研与教学：体育数据分析、基准方法学，以及可复现的预测与推断方法评估。它也可以直接服务速拧社区，例如帮助选手理解自身进步，或帮助组织者推理轮次结构。

\paragraph{禁止用途。}随附数据卡以最强措辞禁止将本数据或模型用于：
\begin{itemize}
  \item 任何形式的\textbf{赌博、博彩或下注}，包括面向投注者的预测服务；
  \item 对个体选手的\textbf{歧视性筛选、画像或羞辱}；
  \item \textbf{冒充} \wca{} 或其他官方机构，或伪造比赛成绩；以及
  \item 未匿名化即二次分发个体级派生数据。
\end{itemize}
这些禁令并非装饰。比赛记录是公开的聚合数据，但其聚合可支持个人未曾同意的推断。预测基准被挪用于博彩的风险足够真实，因此我们把防止它作为设计要求。

\paragraph{署名。}所有派生发表物必须包含 \wca{} 署名声明：“本信息基于 World Cube Association 拥有并维护的比赛结果，发布于 \url{https://worldcubeassociation.org/results}。”代码以 Apache-2.0 发布；数据权利归 \wca{} 所有。

\paragraph{隐私与退出。}尽管底层数据是公开的，个人仍可能有正当理由撤回。因此我们承诺提供退出机制：选手可联系维护者，请求从已发布产物中移除其个体级派生特征。聚合统计（不指向个人）予以保留，使基准的科学主张仍可复现，同时尊重个人选择。我们建议挑战赛组织者在规则中重申禁止用途，并建议 \wca{} Results Team 作为常设联系点。

\paragraph{公平性考量。}由于地区机会不均，主要在代表性较强地区数据上训练的模型，对代表性较弱地区可能表现更差；若被误用，这可能不利于这些地区的选手。我们通过按大洲分层缓解，并拒绝发布个体级排名产品。我们并不声称分层解决了根本的不平等，它只是使之可见。

\paragraph{更广泛影响。}我们希望 \wcabench{} 推动体育机器学习走向更严格的评估：时间感知划分、冻结统计量、分层报告、效应量与显式识别假设。主要社会风险是被用于博彩或个体定向，数据卡与退出机制正是为抑制此类用途而设计。
